% Capítulo 2 — Aplicação empírica

\section{\MakeUppercase{A eólica no Nordeste como captura da transição energética}}

\subsection{Vetor fiscal e acumulação por despossessão}

Aplicada ao recorte cearense e potiguar, a régua de \citeauthoronline{harvey2010} opera com precisão sobre a engrenagem que articula \textit{Regime Especial de Incentivos para o Desenvolvimento da Infraestrutura} (REIDI) à isenção SUDENE.

Com efeito, \citeauthoronline{piketty2014} demonstra que a desigualdade fundamental em que a taxa de retorno do capital supera a taxa de crescimento da economia constitui condição estrutural do capitalismo avançado.

\subsection{Vetor contratual e expulsão pelos contratos eólicos}

Sob esse prisma, \citeauthoronline{ostrom1990} formula sete princípios construtivos para a sustentação de instituições comunais de longa duração. Aplicada ao recorte, a régua de Ostrom diagnostica que os contratos eólicos extinguem o sétimo princípio, ao remover a possibilidade de escolha coletiva entre comunidade anfitriã e operadora.
